%--------------------------
\begin{frame}{Question 2}

\begin{block}{}
L’absence de dépôt d’une demande d’examen préliminaire international :

\begin{itemize}
  \item[$\bullet$] a. n’a aucun effet sur le délai d’ouverture de la phase nationale. 
  \item[$\bullet$] b. ne ralentit pas nécessairement les procédures nationales de délivrance en phase 
  nationale. 
  \item[$\bullet$] c. empêche le déposant de déposer des modifications de la description ou des dessins 
  pour qu’elles soient examinées pendant la phase internationale. 
  \item[$\circ$] d. Aucune des réponses a, b ou c n’est correcte. 
\end{itemize}

\end{block}

\end{frame}

%--------------------------
\begin{frame}{Question 2}

\textbf{Base juridique :} \pctart{22}(1) + \pctart{34}(2)(b).  

\begin{itemize}
  \item[$\bullet$] \textbf{a.} \; Correct : le délai d’entrée en phase nationale
  n’est pas affecté par l’absence de Chapitre II (\pctart{22}(1)).

  \item[$\bullet$] \textbf{b.} \; Correct : l’absence d’examen préliminaire n’implique pas nécessairement
  un ralentissement, car l’examen reste national.

  \item[$\bullet$] \textbf{c.} \; Correct : les modifications de la description et des dessins pendant la phase
  internationale sont examinées uniquement dans le cadre du Chapitre II
  (\pctart{34}(2)(b)).

  \item[$\circ$] \textbf{d.} \; Faux puisque \textbf{a}, \textbf{b} et \textbf{c} sont correctes.
\end{itemize}

\end{frame}
